% =====================================================================
\section{From Safety to Risk: Statistics for Learned Guards}
\label{ch:risk}
% =====================================================================

\begin{bigidea}
So far a guard's contract was \emph{exact}: if the assumptions hold, the promise holds. But a classifier or LLM judge is never exactly right; it errs on some inputs. The paper's answer is not to model each error rate and multiply them (per-layer rates \emph{mislead}). Instead it measures the \textbf{deployed function, route by route}: total risk is the \emph{sum over acceptance routes}; each statistical route gets a \textbf{Clopper--Pearson certificate} (an upper confidence bound from counting losses); a union bound makes all certificates hold together; and a hash of each route's prefix tells you exactly which certificates a change invalidates.
\end{bigidea}

\subsection{Why the exact analysis is not enough}

The composition check of Chapter~\ref{ch:check} treats a guard as if it were right whenever its obligations hold. For a learned guard, an obligation such as \ob{miss\_rate\_bounded} is not a fact about one invocation. It is a claim about a \emph{population}: ``this classifier misses at most such-and-such fraction of attacks''. It stands for an external condition, \emph{a certificate for this configuration, or a bound stated by a source, exists}. The exact analysis tells us \emph{what depends on what} (an idealization: ``if the stage were right whenever its obligations held, would the composition be justified?''). The actual guarantee for statistical goals must come from \textbf{measurement}.

\subsection{Background: just enough probability}

\begin{toolbox}{probability for this chapter}
\textbf{Probability space.} We model proposals $x$ as random draws from a distribution $\mathcal D$ (``the traffic''). An \emph{event} is a set of outcomes, e.g.\ ``$x$ is unsafe and accepted''. $\Pr_{\mathcal D}[E]$ is its probability.

\textbf{Independent and identically distributed (i.i.d.).} A sample of $N$ proposals drawn independently from the same $\mathcal D$. The number $k$ of ``losses'' (unsafe and accepted) among them follows a \textbf{binomial distribution} with parameters $N$ and $p=\Pr[\text{loss}]$.

\textbf{Upper confidence bound.} We never know $p$; we see $k$ out of $N$. A $(1-\delta)$ \emph{upper confidence bound} $u=u(k,N,\delta)$ is a number computed from the sample such that $\Pr[p\le u]\ge 1-\delta$ over the randomness of the sample. For instance $\delta=0.05$ gives a 95\% bound.

\textbf{Union bound (Boole's inequality).} For any events $E_1,\dots,E_m$: $\Pr[E_1\cup\cdots\cup E_m]\le\sum_i\Pr[E_i]$. No independence needed. If each of $m$ statements fails with probability at most $\delta_i$, all hold simultaneously with probability at least $1-\sum_i\delta_i$.

\textbf{Total variation distance.} For two distributions $\mathcal D,\mathcal D'$: $\eta=\sup_E|\mathcal D(E)-\mathcal D'(E)|$ over events $E$, the largest disagreement about any event's probability. If the traffic changes from $\mathcal D$ to $\mathcal D'$ with distance $\eta$, the probability of any event changes by at most $\eta$.
\end{toolbox}

\subsection{The setup: exact goals, statistical goals, joint risk}

Split the goals into two kinds: $\Goals=\Goals_d\uplus\Goals_s$.

\begin{itemize}
\item $\Goals_d$: \textbf{exact goals}, which must hold on \emph{every} accepted proposal (deterministic gates, exact contracts).
\item $\Goals_s$: \textbf{statistical goals}, for which a rate of failure is tolerated (semantic goals that rest on learned guards).
\end{itemize}

Let $Y(x)=1$ if proposal $x$ violates a goal and $0$ otherwise. The \textbf{joint risk} of a deployed configuration (closed routes deny) is
\[
R=\Pr_{\mathcal D}\bigl[Y=1\ \wedge\ x\text{ is accepted}\bigr].
\]
This is the chance that a random proposal is both unsafe \emph{and} let through, the quantity we want small. Note that it is a joint probability. If attacks are rare, $R$ can be tiny even for a useless guard; we return to this below.

\subsection{Risk is a sum over routes}

Recall that every acceptance is on a named route (the Routes property). Group the routes:

\begin{itemize}
\item A route is \textbf{gated} if deterministic gates establish every goal of $\Goals_s$ on every path to it.
\item The other routes form the set $Q$, called \textbf{statistical routes}.
\item A route on which deterministic gates establish \emph{every} goal (exact and statistical) is \textbf{fully gated}; only such routes stand up to an adaptive attacker.
\end{itemize}

\begin{statement}{Theorem (Route risk) \hfill\normalfont paper: Theorem 9}
Let $H$ pass Algorithm~1. Suppose on the support of $\mathcal D$: (i)~deterministic gates decide their goals exactly, in $\Goals_d$ and in $\Goals_s$; and (ii)~realization holds and all contracts hold for goals in $\Goals_d$. Then
\[
R=\sum_{r\in Q}R_r,\qquad R_r=\Pr_{\mathcal D}[Y=1\wedge\accepted(r)].
\]
\end{statement}

\begin{plainproof}{Route risk}
\emph{Step 1: split acceptance into routes.} By the Routes property, an accepted proposal is accepted at exactly one guard. So ``accepted'' is the disjoint union of the events $\accepted(r)$, and $R=\sum_rR_r$ over all routes.

\emph{Step 2: exact goals never fail.} Let $H_d$ be the hierarchy $H$ restricted to $\Goals_d$: replace the goals by $\Goals_d$, intersect each scope with $\Goals_d$, and drop goals outside $\Goals_d$ from gates. Restricting the goals only makes $\ok$ easier at each guard, so $H_d$ passes the check too. By (i) and (ii) the model of $x$ satisfies the contracts of $H_d$, so by Soundness/Exactness an accepted $x$ satisfies $\Goals_d$.

\emph{Step 3: gated routes contribute nothing.} If $x$ is accepted on a gated route, it passed on its path the deterministic gates for every goal of $\Goals_s$, which hold by (i). With Step 2, $x$ violates no goal: $Y(x)=0$, so $R_r=0$ for gated routes. Only the routes in $Q$ remain in the sum. \hfill$\blacksquare$
\end{plainproof}

\textbf{What the theorem is about.} Notice what it does \emph{not} mention: no statistical contract, no independence assumption, no per-layer rate. It says: \emph{the risk of the whole system is the sum of route risks, and routes that deterministic gates have fully covered contribute zero.} The composition check thus tells you which routes ($Q$) need statistical evidence.

\begin{example}{DevAgent routes}
With $\Goals_s=\{\varphi_3,\varphi_4\}$: route $n_3$ (the judge) has $\varphi_3$ established by gate $c$ but $\varphi_4$ is covered only by the learned judge, so it is statistical. Route $n_4$ (the human) covers both. So $Q=\{n_3,n_4\}$, and $R=R_{n_3}+R_{n_4}$.
\end{example}

\subsection{Per-layer rates mislead}

A frequent informal argument says: ``Layer 1 misses 10\% of attacks, layer 2 misses 20\%, layer 3 misses 30\%. All three must miss for an attack to succeed, so the system misses $0.1\times0.2\times0.3=0.6\%$.'' This multiplies marginal rates, which is correct \emph{only if the layers' misses are independent}. They are typically not: an attack that fools one classifier tends to fool similar ones.

\begin{statement}{Proposition (Veto and cascade) \hfill\normalfont paper: Proposition 2}
Let guards have marginal miss probabilities $m_i$ on the unsafe population.\\
(a)~In a \emph{veto chain}, the joint miss probability can be anything from $\max\bigl(0,\,1-\sum_i(1-m_i)\bigr)$ to $\min_im_i$, and $\prod_im_i$ is one value in that range, attained under conditional independence.\\
(b)~In a \emph{cascade with early acceptance}, neither $\min_im_i$ nor $\prod_im_i$ bounds risk: if a first guard accepts a share $\rho$ of unsafe proposals and a later guard blocks all it sees, the later guard's miss probability is zero, so the minimum and the product are zero, while the risk is $\rho$ times the unsafe mass.
\end{statement}

\begin{plainproof}{Veto and cascade}
\textbf{(a)} These are the \textbf{Fr\'echet--Hoeffding bounds}. The intuition: ``missed by guard $i$'' is an event $M_i$ with $\Pr[M_i]=m_i$. The joint event $\bigcap_iM_i$ is smallest if the misses are spread out as much as possible (no overlap beyond what is forced), and largest if the misses are nested ($M_i$ contained in the smaller one). The upper bound $\min_im_i$ is attained when the misses are nested: $\Pr[\bigcap M_i]\le\Pr[M_i]$ for each $i$. The lower bound comes from the union bound applied to the complements: $\Pr[\bigcap M_i]=1-\Pr[\bigcup M_i^c]\ge1-\sum_i(1-m_i)$.

\textbf{(b)} By construction and the Routes property: unsafe proposals accepted by the first guard never meet the later one; the later guard's miss probability is measured on the population it \emph{sees}. \hfill$\blacksquare$
\end{plainproof}

\begin{example}{the numbers $0.1$, $0.2$, $0.3$}
With $m=(0.1,0.2,0.3)$: lower bound $\max(0,1-(0.9+0.8+0.7))=\max(0,-1.4)=0$; upper bound $\min=0.1$; independence gives $0.006$. All values in $[0,0.1]$ are possible. The repository's lab computes joint risk under a ``common-shock'' knob that interpolates between independence and nesting: weight $0\to0.006$, $0.5\to0.053$, $1\to0.100$. A factor of about $17$ separates the optimistic and pessimistic readings of the \emph{same} three miss rates.
\end{example}

\begin{tryit}{see the spread yourself}
# from the repository root
python scripts/tutorial_lab.py
# part 1: same guard miss rates, different joint risks (0.006 ... 0.100)
# part 2: pairwise independence does not determine triple risk
# part 3: zero observed errors still leave uncertainty
\end{tryit}

\begin{example}{why the cascade product is wrong}
Suppose 10\% of the unsafe proposals are accepted by an early, weak guard $G_1$ (it escalates the other 90\%). A later guard $G_2$ blocks everything it sees. Then $G_2$'s miss rate (on what it sees) is $0$, so $\min_im_i=0$ and $\prod_im_i=0$, and a naive calculation announces ``risk 0''. But 10\% of the unsafe mass escapes via $G_1$. The route-sum picture gets it right: $R_{G_1}=0.10\times(\text{unsafe mass})$.
\end{example}

\subsection{Route certificates}

\textbf{The plan.} Do not model the guards. Instead, for each statistical route $r$, \emph{run the deployed decision function} on $N$ labelled samples, and count how many were unsafe and accepted at that route. Turn the count into an upper confidence bound.

\begin{defbox}{route certificate}
A tuple $(r,f,c,u_r,\delta_r,N)$: the route $r$; a \textbf{prefix fingerprint} $f$ (a hash of everything upstream of $r$, see below); a context label $c$ naming the deployment; an upper bound $u_r$; the confidence share $\delta_r$; and the sample count $N$.
\end{defbox}

The confidence budget $\delta$ (e.g.\ $0.05$) is split into shares $\delta_r$ over the routes of $Q$ \emph{before} any calibration: $\sum_{r\in Q}\delta_r\le\delta$.

\subsubsection{The Clopper--Pearson bound}

If $k_r$ is the number of unsafe proposals accepted on $r$ among $N$ samples, then
\[
u_r=\mathrm{Beta}^{-1}\bigl(1-\delta_r;\ k_r+1,\ N-k_r\bigr),\qquad u_r=1\text{ if }k_r=N.
\]
The inverse Beta function looks intimidating. Here is the idea behind it, and a special case that you can compute by hand.

\begin{toolbox}{the Clopper--Pearson idea}
Suppose the true loss probability were $p$. Then the number of losses among $N$ samples is $\mathrm{Bin}(N,p)$. We observed $k$ losses. \textbf{The bound $u$ is the largest $p$ that is still ``plausible''}, precisely: the value of $p$ at which the probability of seeing $k$ \emph{or fewer} losses equals $\delta$:
\[
\Pr[\mathrm{Bin}(N,u)\le k]=\delta.
\]
For any true $p>u$, seeing $\le k$ losses would have probability less than $\delta$. So with probability at least $1-\delta$ (over the sample), the true $p$ is at most $u$. The Beta-function formula is simply a closed form for this equation.

\smallskip
\textbf{The case $k=0$ by hand.} $\Pr[\mathrm{Bin}(N,u)=0]=(1-u)^N=\delta$ gives $u=1-\delta^{1/N}$. For $N=1800$ and $\delta=0.025$: $\ln\delta=-3.689$, so $\delta^{1/1800}=e^{-0.00205}\approx0.99795$ and $u\approx0.00205$. This matches the paper's value $0.0020$. For small $u$, $u\approx\ln(1/\delta)/N$. (This is the ``rule of three'' for $\delta=0.05$: $u\approx3/N$.) The other number in the paper, ten losses, gives $0.0102$; zero observed losses still leaves an upper bound above zero.
\end{toolbox}

\begin{example}{DevAgent certificate}
$\delta=0.05$ split evenly across the two routes: $\delta_r=0.025$. Take $N=1800$ labelled samples. If ten unsafe proposals were accepted on route $n_3$, then $u_{n_3}=0.0102$: with 97.5\% confidence, the chance that a random proposal is unsafe \emph{and} accepted by the judge is at most about 1\%. Route $n_4$ (the person) is certified the same way from decisions recorded in shadow mode, so that the person's error rate is measured instead of assumed.
\end{example}

\begin{tryit}{compute a bound}
import sys; sys.path[:0] = ["src"]
from tint.calibration import binomial_upper
print(binomial_upper(0, 1800, 0.025))    # about 0.0020 (zero losses)
print(binomial_upper(10, 1800, 0.025))   # about 0.0102 (ten losses)
\end{tryit}

\subsubsection{The runtime rules}

At runtime a statistical route accepts only while a certificate for its \emph{current} fingerprint is open. The engine refuses (i) an allocation above $\delta$, (ii) a certificate above its route's share, (iii) two certificates for one route, and (iv) bounds that sum above the risk budget $\varepsilon$. A \textbf{closed} route denies. So a route with a stale or missing certificate fails \emph{closed}: its traffic is denied (or, as a policy choice, sent to a person).

\subsubsection{The certified-risk theorem}

\begin{statement}{Theorem (Certified risk) \hfill\normalfont paper: Theorem 10}
Under the premises of Theorem~9 (Route risk), fix a configuration whose stages on statistical routes do not read session state. Let shares $\delta_r$ be fixed before calibration with $\sum_{r\in Q}\delta_r\le\delta$, and let each certificate's samples be i.i.d.\ from $\mathcal D$, drawn after its route's prefix was fixed. With probability at least $1-\delta$,
\[
R\ \le\ \sum_{r\ \text{open}}u_r.
\]
The runtime keeps this sum at most $\varepsilon$.
\end{statement}

\begin{plainproof}{Certified risk}
Each $u_r$ is an upper bound on $R_r$ at level $\delta_r$ (the Clopper--Pearson guarantee: $\Pr[R_r>u_r]\le\delta_r$). By the union bound, the probability that \emph{some} route has $R_r>u_r$ is at most $\sum_{r\in Q}\delta_r\le\delta$ \emph{without any independence among the certificates}, so they may come from different runs. Because the shares cover $Q$ and are fixed in advance, which certificates end up open cannot influence the bound. With probability $\ge1-\delta$, then, $R_r\le u_r$ for all $r\in Q$. Now: $R=\sum_{r\in Q}R_r$ (Route risk). Closing a route only removes acceptances (Monotonicity), so closed routes contribute zero; open ones contribute at most $u_r$. Gated routes contribute zero by the Route risk theorem. Therefore $R\le\sum_{r\text{ open}}u_r$. \hfill$\blacksquare$
\end{plainproof}

\subsubsection{Reading the fine print}

\begin{warning}{Four ways to misuse a certificate}
\begin{enumerate}
\item \textbf{Selection.} The statement is per configuration. If you recalibrate repeatedly until some bound passes, you have run many tests and reported the winner; that voids the level unless every attempt is charged to $\delta$. Likewise, over a long-lived deployment with many versions, allocate shares across versions in advance, and reserve shares for routes you expect to add (the engine refuses old certificates whose shares no longer fit).
\item \textbf{Rare attacks make joint bounds look good.} With prevalence $0.001$ of attacks, a guard that misses \emph{every} attack loses only about $2$ of $1800$ samples, and its bound $0.0040$ passes a budget of $0.01$. So also report the \emph{conditional} miss rate ($R$ divided by the unsafe mass), by attack class, and use a calibration set with enough attacks.
\item \textbf{Drift.} A certificate is about the distribution $\mathcal D$ it was calibrated on. If deployment traffic is $\mathcal D'$ at total-variation distance $\eta$, the risk can be higher by at most $\eta$ (risk is the probability of an event), and monitoring should track it.
\item \textbf{Labels.} Where the labels $Y$ come from is usually the hardest part. Nothing in the paper supplies them.
\end{enumerate}
\end{warning}

\subsection{Incremental recertification}

\textbf{Motivation.} Engineers change stages all the time: retrain the scorer, tweak a rule, update the judge's prompt. After each change, which certificates are still valid? Recomputing all of them wastes labelled data and time.

\begin{toolbox}{hash functions and Merkle-style fingerprints}
A \textbf{cryptographic hash function} maps any input to a short fixed-size string (a \emph{digest}); a tiny change to the input changes the digest completely, and it is practically impossible to find two different inputs with the same digest. A \textbf{Merkle} construction hashes a structure by hashing each part together with the digests of the parts it depends on. If anything upstream changes, the digest of everything downstream changes.
\end{toolbox}

Each route is bound to a \textbf{prefix fingerprint}, a Merkle hash computed in topological order. A node's fingerprint covers: its stage and layer; the identities of its rule, guard or condition and of its checked predicates; obligation kinds; the labelled fingerprints of its predecessors; and a digest of the goals, the statistical goals and a label naming the oracle $Y$. A component's \emph{identity} covers its source code, closure values, defaults, partial arguments, reachable instance state, array contents and declared version and configuration. A component with no readable identity is refused. What is \emph{not} covered (code reached only through calls, state outside the object, model-weight files) the deployer must declare in the configuration.

\begin{center}
\begin{tikzpicture}[font=\scriptsize,
  nd/.style={draw=tintblue,rounded corners=2pt,fill=tintblue!8,align=center,minimum height=0.7cm,minimum width=1.3cm},
  st/.style={draw=tintred,rounded corners=2pt,fill=tintred!12,align=center,minimum height=0.7cm,minimum width=1.3cm},
  arr/.style={-{Stealth[length=1.8mm]}}]
\node[nd] (g) at (0,0) {$g_1..g_3,s$};
\node[nd] (v) at (-3,-1.4) {$v_2$}; \node[nd] (h) at (-3,-2.8) {$n_4$ route};
\node[nd] (c) at (3,-1.4) {$c$}; \node[nd] (j) at (3,-2.8) {$n_3$ route};
\draw[arr] (g)--(v); \draw[arr] (v)--(h); \draw[arr] (g)--(c); \draw[arr] (c)--(j);
\node[st,right=1.0cm of v] (vv) {retrain $v_2$};
\node[font=\scriptsize,below=0.1cm of h,text=tintred] {fingerprint changes: certificate closed};
\node[font=\scriptsize,below=0.1cm of j,text=tintgreen!60!black] {fingerprint unchanged: certificate still open};
\end{tikzpicture}
\end{center}

By Corollary 2 (Affected routes), a change to a stage cannot alter the event of a route it does not precede, so only routes it precedes need new evidence. They close alone, and the others keep their certificates. Two caveats from the paper: the cost is a tighter guarantee foregone (for few routes, one whole-hierarchy bound beats a sum of route bounds at stricter levels); and closing a route is an outage for its traffic, so calibrate the new prefix in \emph{shadow mode} before switching. A change that only tightens a stage cannot raise $R_r$ on its own route (Monotonicity), but it can move traffic to later routes, so the conservative rule is kept.

\begin{example}{what a change stales in DevAgent}
\begin{center}
\renewcommand{\arraystretch}{1.25}
\begin{tabular}{@{}ll@{}}
\toprule
\textbf{Changed stage(s)} & \textbf{Certificates closed}\\
\midrule
$g_1,g_2,g_3,s$ (rules, taint switch) & both routes ($n_3$ and $n_4$)\\
$v_2$ or $n_4$ (scorer, human) & route $n_4$ only\\
$c$ or $n_3$ (provenance gate, judge) & route $n_3$ only\\
\bottomrule
\end{tabular}
\end{center}
So a retrained scorer costs one route; a rule change costs all of them. Rule changes may be the most frequent, and for those the incremental scheme gains nothing.
\end{example}

\pauseq{A team replaces the judge $n_3$ by a newer model. Which certificates are stale? Now suppose instead the team adds a new deterministic gate right after $g_3$. What happens?}{Only the certificate of route $n_3$ is stale in the first case (the judge precedes no other route; the route $n_4$'s fingerprint is built from $g_1,g_2,g_3,s,v_2,n_4$ and does not include $n_3$). In the second case, the new gate is an ancestor of both routes, so both fingerprints change and both certificates close until recalibrated.}

\begin{recap}
\begin{itemize}
\item Exact analysis tells \emph{what depends on what}; statistics supplies numbers for learned stages.
\item $R=\sum_{r\in Q}R_r$: risk is a sum over statistical routes; gated routes contribute zero. No independence assumption is needed.
\item Per-layer rates mislead: Fr\'echet--Hoeffding says the joint miss rate can be anywhere between $\max(0,1-\sum(1-m_i))$ and $\min m_i$; for cascades with early acceptance even those bounds fail.
\item Clopper--Pearson: the largest loss rate consistent with $k$ losses in $N$ samples at level $\delta_r$. Union bound makes all route certificates hold together.
\item Fingerprints (Merkle hashes of the prefix) tell exactly which certificates a change invalidates; they are only as good as the declared component identity.
\end{itemize}
\end{recap}
